\subsection{A polynomial modular-index bound for coefficient positivity}

\begin{proposition}[P-F11: positive rational form implies a polynomial positivity-index bound]
\label{prop:modularbound}
Let $f(n)$, $n\ge0$, be a quasi-polynomial with nonzero generating function
\[
 F(T)=\frac{H(T)}{\prod_{i=1}^r(1-T^{a_i})},\qquad
 H\in\mathbb Z_{\ge0}[T]\setminus\{0\},\quad a_i>0.
\]
Here $r\ge1$. Set $A=\max_i a_i$ and $h=\deg H$. Then
\[
 p(f)\le h+128r^3A^3 < rA+128r^3A^3.
 \tag{P-F11}
\]
The bound concerns the positivity index of every constituent polynomial,
not just eventual positivity of the values in its own congruence class.
In particular, polynomial bounds for $r$ and the modular index $A$
suffice for a polynomial bound for the positivity index; a polynomial
bound on the common period is unnecessary.
\end{proposition}
\begin{proof}
The inequality $h<\sum_i a_i\le rA$ was proved in P-F9 from global
quasi-polynomiality. It remains to prove the first bound. Begin with
\[
 G(T)=\prod_{i=1}^r(1-T^{a_i})^{-1},\qquad g(n)=[T^n]G(T),
\]
assuming $\gcd(a_1,\ldots,a_r)=1$ and $r\ge2$. Put
$D=16rA^2$ and $\rho=1/D$. Every pole $\zeta$ of $G$ is a root of unity
of order at most $A$. Its order as a pole is
$m=m(\zeta)=\#\{i:\zeta^{a_i}=1\}$.
There are at most $N:=\sum_i a_i\le rA$ distinct poles. The only pole
of order $r$ is $\zeta=1$, by the greatest-common-divisor assumption.
Write the partial fractions as
\[
 G(T)=\sum_{\zeta}\sum_{j=1}^{m(\zeta)}
 C_{\zeta,j}(1-T/\zeta)^{-j}.
\]
There is no polynomial part because $G$ is proper, and
$C_{1,r}=c:=1/\prod_i a_i>0$.

We first establish, for every displayed term,
\[
 |C_{\zeta,j}|/c\le 2D^{r-j}.                 \tag{*}
\]
Use the local variable $z=1-T/\zeta$ and put
$u_a(z)=(1-(1-z)^a)/(az)=a^{-1}\sum_{k=0}^{a-1}(1-z)^k$.
For $|z|\le\rho$,
\[
 |u_a(z)-1|\le\tfrac12(a-1)\rho e^{a\rho}\le a\rho,
 \qquad
 |(1-z)^a-1|\le a\rho e^{a\rho}\le2a\rho.
\]
These inequalities use $a\rho\le1/16$ and the finite geometric sum;
the last constants are deliberately loose. If $\zeta$ has order $d$
and $\zeta^{a_i}\ne1$, then
\[
 |1-\zeta^{a_i}|\ge2\sin(\pi/d)\ge4/d\ge4/A.
\]
Consequently
$|1-\zeta^{a_i}(1-z)^{a_i}|\ge4/A-2A\rho>3/A$.
For the $m$ factors vanishing at $z=0$, the product of the inverses of
$u_{a_i}$ is at most
\[
 \prod_{\zeta^{a_i}=1}(1-a_i\rho)^{-1}
 \le\exp(2\rho\textstyle\sum_i a_i)
 \le\exp(1/(8A))<2.
\]
Thus the analytic function $z^mG(\zeta(1-z))/c$ has absolute value at
most $2A^{2(r-m)}$ on $|z|\le\rho$: each of the other factors contributes
at most $a_i/(3/A)\le A^2$.
Cauchy's coefficient bound gives
\[
 |C_{\zeta,j}|/c
 \le2A^{2(r-m)}\rho^{-(m-j)}
 \le2D^{r-j},
\]
which is (*). The denominator estimates also show explicitly that this
analytic function has no further pole in that disk.

Choose a period divisible by the orders of all poles. On a fixed residue
class $b$, partial fractions give the constituent polynomial
\[
 g_b(n)=c\binom{n+r-1}{r-1}
 +\sum_{(\zeta,j)\ne(1,r)}
     C_{\zeta,j}\zeta^{-b}\binom{n+j-1}{j-1}.
\]
This identity holds for all nonnegative integers on the class, hence as
a polynomial identity for the constituent. Let $p>0$ be real and let
$[n^k]$ denote a polynomial coefficient. All coefficients of
$B_j(n):=\binom{n+p+j-1}{j-1}$ are nonnegative. For $j<r$ and
$0\le k\le j-1$,
\[
 \frac{[n^k]B_j(n)}{[n^k]B_r(n)}
 \le \frac{(r-1)!}{(j-1)!p^{r-j}}
 \le (r/p)^{r-j}.                              \tag{**}
\]
Indeed, multiply the $j-1$ factors in $B_j$ by the additional factors
$(n+p+j),\ldots,(n+p+r-1)$, whose constant product is at least
$p^{r-j}$, and compare coefficients. For $k>j-1$ the numerator is zero.

Using (*) and (**), the absolute value of the total error coefficient,
divided by the corresponding positive coefficient of $cB_r$, is at most
\[
 2N\sum_{\ell=1}^{r-1}(Dr/p)^\ell.
\]
At $p_0=128r^3A^3$ one has $Dr/p_0=1/(8rA)\le1/(8N)$.
The last display is bounded by
$2N(1/(8N))/(1-1/(8N))\le2/7<1$.
Each constituent has real coefficients, since it agrees with integer
values at infinitely many real integers. The strict bound on the
absolute value of its complex error therefore shows that every
coefficient of every $g_b(n+p_0)$ is strictly positive.
The same holds for every real $p\ge p_0$, either by the same estimates
or by a further nonnegative translation.

Now let $d=\gcd(a_1,\ldots,a_r)$ be arbitrary and set $a_i'=a_i/d$.
The coefficients of $G$ vanish outside multiples of $d$, and on those
multiples equal $g'(n/d)$ for the relatively prime weights $a_i'$.
A nonzero constituent is correspondingly a constituent of $g'$ evaluated
at $n/d$. It has nonnegative coefficients after every shift
$p\ge d\cdot128r^3(A/d)^3\le128r^3A^3$; zero constituents remain zero.
If $r=1$, the coefficient sequence is simply the indicator of divisibility
by $a_1$, and already has nonnegative constant constituents.

Finally, $F(T)=\sum_i H_iT^iG(T)$ with $H_i\ge0$.
On sufficiently large indices its constituent polynomials are the sums
of $H_i$ times the appropriately relabeled constituents of $g(n-i)$.
The same polynomial identities hold everywhere because $f$ is globally
a quasi-polynomial. After shifting $n$ by $h+p_0$, each summand has
shift $h+p_0-i\ge p_0$, so every coefficient is nonnegative.
This proves the asserted positivity-index bound and completes the proof.
\end{proof}

P-F11 is an elementary analytic bound on rational generating functions;
its complex analysis is Cauchy's coefficient estimate on a fixed small
disk, and it uses neither RH nor a purity theorem. It strengthens the
common-period sufficient condition P-F9. A historical novelty assertion
would require a separate literature review; the statement is included
here as a fully proved deduction, not as a claim of priority.
